% preamble/cover.tex
% 简化版封面 —— Dummit 抽象代数读书笔记
% 蓝金配色，保留九层 TikZ 设计骨架

\makeatletter
\providecommand{\@coverimage}{}
\newcommand{\coverimage}[1]{\renewcommand{\@coverimage}{#1}}

\newcommand{\makecover}{%
\begin{titlepage}
\thispagestyle{empty}

% ============================================================
%  第一层：左侧深蓝色竖条 + 装饰
% ============================================================
\begin{tikzpicture}[remember picture,overlay]
  \fill[RedDark!95]
    (current page.south west) rectangle
    ([xshift=2.8cm]current page.north west);
  \shade[top color=RedDark!90, bottom color=RedDark!95]
    ([xshift=0cm,yshift=-4cm]current page.north west) rectangle
    ([xshift=2.8cm]current page.north west);
  \shade[bottom color=RedDark!90, top color=RedDark!95]
    ([xshift=0cm,yshift=4cm]current page.south west) rectangle
    ([xshift=2.8cm]current page.south west);
  \draw[RedAccent!50,line width=1pt]
    ([xshift=2.8cm]current page.south west) --
    ([xshift=2.8cm]current page.north west);
  \draw[RedLight!25,line width=0.4pt]
    ([xshift=0.4cm]current page.south west) --
    ([xshift=0.4cm]current page.north west);
  \foreach \y/\r/\o in {
    0.08/1.1/0.07, 0.22/0.7/0.11, 0.40/1.4/0.06,
    0.55/0.55/0.10, 0.68/0.9/0.08, 0.82/0.6/0.11, 0.93/0.75/0.09
  } {
    \fill[white, opacity=\o]
      ([xshift=1.4cm, yshift=-\y\paperheight]current page.north west)
      circle (\r cm);
  }
  \draw[white, line width=0.3pt, opacity=0.18]
    ([xshift=0.5cm]current page.north west) --
    ([xshift=0.5cm]current page.south west);
  \draw[white, line width=0.3pt, opacity=0.18]
    ([xshift=2.1cm]current page.north west) --
    ([xshift=2.1cm]current page.south west);
\end{tikzpicture}

% ============================================================
%  第二层：顶部/底部装饰金线
% ============================================================
\begin{tikzpicture}[remember picture,overlay]
  \draw[GoldMain!50,line width=0.6pt]
    ([xshift=3.0cm,yshift=-1.6cm]current page.north west) --
    ([xshift=-2.0cm,yshift=-1.6cm]current page.north east);
  \foreach \x in {3.5,7,10.5,14} {
    \fill[GoldMain!60]
      ([xshift=\x cm,yshift=-1.45cm]current page.north west)
      circle (1.5pt);
  }
  \draw[GoldMain!40,line width=0.5pt]
    ([xshift=3.0cm,yshift=2.2cm]current page.south west) --
    ([xshift=-1.5cm,yshift=2.2cm]current page.south east);
  \foreach \x in {4,8,12} {
    \fill[GoldMain!45]
      ([xshift=\x cm,yshift=2.35cm]current page.south west)
      circle (1.2pt);
  }
\end{tikzpicture}

% ============================================================
%  第三层：右上角装饰几何图形
% ============================================================
\begin{tikzpicture}[remember picture,overlay]
  \draw[GoldMain!30,line width=0.7pt]
    ([xshift=-2.5cm,yshift=-2.5cm]current page.north east)
    circle (22pt);
  \draw[RedAccent!30,line width=0.5pt]
    ([xshift=-2.5cm,yshift=-2.5cm]current page.north east)
    circle (14pt);
  \fill[GoldMain!50]
    ([xshift=-2.5cm,yshift=-2.5cm]current page.north east)
    circle (2.5pt);
  \foreach \ang in {0,90,180,270} {
    \draw[GoldMain!25,line width=0.4pt]
      ([xshift=-2.5cm,yshift=-2.5cm]current page.north east)
      -- ++(\ang:28pt);
  }
  \foreach \ang in {45,135,225,315} {
    \draw[RedAccent!20,line width=0.3pt]
      ([xshift=-2.5cm,yshift=-2.5cm]current page.north east)
      -- ++(\ang:26pt);
  }
\end{tikzpicture}

% ============================================================
%  第四层：左下角装饰曲线
% ============================================================
\begin{tikzpicture}[remember picture,overlay]
  \draw[GoldMain!15,line width=0.5pt,smooth]
    ([xshift=3.5cm,yshift=3.0cm]current page.south west)
    to[out=30,in=210] ++(3.0,2.0);
  \draw[RedAccent!15,line width=0.5pt,smooth]
    ([xshift=3.3cm,yshift=3.0cm]current page.south west)
    to[out=35,in=205] ++(3.2,2.2);
  \draw[RedMain!12,line width=0.5pt,smooth]
    ([xshift=3.5cm,yshift=4.0cm]current page.south west)
    to[out=-15,in=195] ++(2.5,-1.5);
\end{tikzpicture}

% ============================================================
%  第五层：公式背景区
% ============================================================
\begin{tikzpicture}[remember picture,overlay]
  \shade[top color=GoldPale!40, bottom color=GoldPale!25]
    ([xshift=3.3cm,yshift=-3cm]current page.north west) rectangle
    ([xshift=-1.5cm,yshift=2.5cm]current page.south east);
  \draw[GoldMain!30,line width=0.4pt]
    ([xshift=3.3cm,yshift=-3cm]current page.north west) rectangle
    ([xshift=-1.5cm,yshift=2.5cm]current page.south east);
  
  % 代数学经典公式
  \node[anchor=west,RedMain!70] at ([xshift=3.7cm,yshift=-3.5cm]current page.north west)
    {$\displaystyle G/\Ker\phi\;\cong\;\Im\phi$};
  \node[anchor=west,RedAccent!65] at ([xshift=3.8cm,yshift=-4.8cm]current page.north west)
    {$\displaystyle |G|=|H|\cdot[G:H]$};
  \node[anchor=west,GoldMain!75] at ([xshift=3.9cm,yshift=-6.1cm]current page.north west)
    {$\displaystyle \Aut(\Z/n\Z)\cong(\Z/n\Z)^{\times}$};
  \node[anchor=west,RedMain!55] at ([xshift=4.0cm,yshift=-7.4cm]current page.north west)
    {$\displaystyle \Gal(E/F)\cong G$};
  \node[anchor=west,RedAccent!50] at ([xshift=3.7cm,yshift=-8.7cm]current page.north west)
    {$\displaystyle R/I\ \text{域}\ \Longleftrightarrow\ I\ \text{极大理想}$};
  
  \foreach \y/\c in {-3.5/RedMain!35, -4.8/RedAccent!30, -6.1/GoldMain!35, -7.4/RedMain!30, -8.7/RedAccent!25} {
    \fill[\c]
      ([xshift=-2.0cm,yshift=\y cm]current page.north west)
      circle (1.8pt);
  }
\end{tikzpicture}

% ============================================================
%  第六层：标题区
% ============================================================
\vspace*{2cm}

\hspace*{3.3cm}
\begin{minipage}{0.7\textwidth}
  {\large\bfseries\color{RedAccent} Dummit \& Foote 抽象代数\quad 读书笔记\par}
  
  \vspace{0.3cm}
  \begin{tikzpicture}
    \draw[GoldMain!60,line width=1.2pt] (0,0) -- (2.0,0);
    \fill[GoldMain!70] (2.3,0) circle (2pt);
    \fill[GoldMain!50] (2.8,0) circle (1.2pt);
    \fill[GoldMain!30] (3.2,0) circle (0.8pt);
  \end{tikzpicture}
  
  \vspace{0.8cm}
  {\Huge\bfseries\color{RedDark} \@title\par}
  
  \vspace{0.6cm}
  {\Large\color{RedMain!70} \@author\par}
  
  \vspace{0.4cm}
  \begin{tikzpicture}
    \draw[GoldMain!40,line width=0.7pt] (0,0) -- (3.5,0);
    \fill[GoldMain!45] (0,0) circle (1.5pt);
    \fill[GoldMain!45] (3.5,0) circle (1.5pt);
  \end{tikzpicture}
  
  \vspace{0.8cm}
  {\large\color{GoldMain!85} 译注：万一之\par}
  \vspace{0.5cm}
  {\normalsize\color{RedDark!60} 原书：\textit{Abstract Algebra}, 3rd ed.\par}

\end{minipage}

% ============================================================
%  第七层：底部图片
% ============================================================
\ifdefempty{\@coverimage}{}{
	\begin{tikzpicture}[remember picture,overlay]
		\node[anchor=south west,inner sep=0pt,outer sep=0pt]
		at ([xshift=3.3cm,yshift=3.05cm]current page.south west)
		{
			\includegraphics[
			width=\dimexpr\paperwidth-4.8cm\relax
			]{\@coverimage}
		};
	\end{tikzpicture}
}

% ============================================================
%  第八层：散落装饰点
% ============================================================
\begin{tikzpicture}[remember picture,overlay]
  \fill[GoldMain!20] ([xshift=-3.0cm,yshift=-9.5cm]current page.north west) circle (1.0pt);
  \fill[RedAccent!20] ([xshift=-3.8cm,yshift=-9.8cm]current page.north west) circle (0.7pt);
  \fill[RedMain!15] ([xshift=-4.0cm,yshift=3.5cm]current page.south west) circle (1.2pt);
  \fill[GoldMain!15] ([xshift=-5.5cm,yshift=4.0cm]current page.south west) circle (0.9pt);
\end{tikzpicture}

\end{titlepage}
}
\makeatother
